\begin{pastexam}[label={px:ncu-110-dsa-12}]{中央資工所 110 學年度 問答題第 2 題}
（問答）

A weighted graph G=(V, E) is an undirected graph with nonnegative weight assigned to each edge in
E. The length of a path in G is the sum of edge weights of the edges in the path. A path P from
vertex u to vertex v is said to be shortest if the length of P is smallest for all paths from vertex
u to vertex v.
\begin{enumerate}[label=(\Alph*)]
\item (12\%) For a weighted graph $G$=($V$, $E$) and two vertices $x$, $y$ in $V$, give an efficient
algorithm to find a shortest path from $x$ to $y$. Analyses the time complexity of the algorithm.
\item (13\%) Now you are given a weighted graph $G$=($V$, $E$), a specified vertex $s$ in $V$, an
array $d$[\,] of size $|V|$, and it is claimed that for each $v$ in $V$, $d$[$v$] is the length of the
shortest path from $s$ to $v$. The claim may be correct or may be not. Design a linear time
algorithm to check whether the claim is correct. Analyses the time complexity of your algorithm and
make sure it is linear, that is, the running time of your algorithm must be $O(|V|+|E|)$.
\end{enumerate}
\end{pastexam}

\begin{pastsoln}
記 $n = |V|$、$m = |E|$，$\delta(u, v)$ 是 $u$ 到 $v$ 真正的最短距離，走不到時 $\delta = \infty$。 圖用
鄰接串列存，無向邊 $\{u, v\}$ 在 $u$、$v$ 的串列裡各出現一次（看成兩條方向相反、權重相同的有向邊）。
權重都 $\ge 0$，所以沒有負環，最短路徑一定存在（走得到的話）。

\textbf{(A)} 權重非負，用 Dijkstra 演算法（資結書第~25.2~節、虛擬碼~25.3），以 $x$ 為起點；$y$ 一被
取出就可以停，再沿著前驅 $\pi$ 從 $y$ 走回 $x$。
\begin{pseudo}{$x$ 到 $y$ 的最短路徑}
\Function{Shortest-Path}{$G, w, x, y$}
    \For{每個頂點 $v \in V$}
        \State $d[v] \gets \infty$；$\pi[v] \gets \NIL$；$\mathit{done}[v] \gets \FALSE$
    \EndFor
    \State $d[x] \gets 0$
    \State $Q \gets$ 以 $d$ 為鍵、含所有頂點的最小優先佇列
    \While{$Q$ 不是空的}
        \State $u \gets$ \Call{Extract-Min}{$Q$}
        \If{$d[u] = \infty$ \OR $u = y$}
            \State \Break \Comment{剩下的都走不到，或 $y$ 已經確定}
        \EndIf
        \State $\mathit{done}[u] \gets \TRUE$
        \For{每個 $v \in \mathit{Adj}[u]$}
            \If{$\mathit{done}[v] = \FALSE$ \AND $d[u] + w(u, v) < d[v]$}
                \State $d[v] \gets d[u] + w(u, v)$；$\pi[v] \gets u$
                \State \Call{Decrease-Key}{$Q, v, d[v]$}
            \EndIf
        \EndFor
    \EndWhile
    \If{$d[y] = \infty$}
        \State \Return 「$x$ 走不到 $y$」
    \EndIf
    \State 從 $y$ 沿著 $\pi$ 走回 $x$，把經過的頂點倒過來，\Return 這條路徑和長度 $d[y]$
\EndFunction
\end{pseudo}

\textbf{正確性}。 資結書定理~25.6：權重非負時，每個頂點被 \proc{Extract-Min} 取出的那一刻，$d[u] =
\delta(x, u)$，而且之後不會再變。 所以 $y$ 被取出時 $d[y] = \delta(x, y)$，可以提早結束。 取出的頂點
$d = \infty$ 時，佇列裡剩下的頂點 $d$ 也都是 $\infty$，它們（包括還沒取出的 $y$）都從 $x$ 走不到。
路徑的部分：$\pi[v]$ 最後一次被設成 $u$ 時，$u$ 已經取出，$d[u] = \delta(x, u)$ 不再改變，而且
$d[v] = d[u] + w(u, v)$。 從 $y$ 沿著 $\pi$ 往回走，經過的都是已經取出的頂點，每一步都滿足這個等式，
一路加回 $d[x] = 0$，得到一條長度恰好是 $d[y] = \delta(x, y)$ 的路徑，也就是最短路徑（和資結書
定理~25.5 同一個道理）。

\textbf{時間}。 初始化 $O(n)$。 每個頂點最多被取出一次，每條無向邊在兩個端點的串列裡各被檢查最多一次，
所以最多 $n$ 次 \proc{Extract-Min}、$2m$ 次 \proc{Decrease-Key}。 取回路徑最多 $n - 1$ 步。 依優先佇列
的做法：
\begin{itemize}
\item 二元堆積：每個操作 $O(\log n)$，總共 $O((n + m)\log n)$。 圖是連通的時候 $m \ge n - 1$，就是
$O(m \log n)$。
\item 陣列（每次掃一遍找最小）：\proc{Extract-Min} $O(n)$、\proc{Decrease-Key} $O(1)$，共 $O(n^2 + m) =
O(n^2)$，適合稠密圖。
\item 費氏堆積：\proc{Decrease-Key} 平攤 $O(1)$，共 $O(m + n \log n)$（資結書第~8.6~節）。
\end{itemize}
提早結束不會讓最壞情況變好（$y$ 可能是最後一個被取出的），但平均上常常省很多。 權重全部相同時可以
直接用 BFS，$O(n + m)$；這題權重任意非負，所以用 Dijkstra。

\textbf{(B)} 要同時確認兩件事：$d[v]$ 不比真正的距離大（真的有一條那麼短的路），也不比它小（沒有
更短的路）。 「沒有更短的路」用三角不等式檢查每條邊；「真的有那麼短的路」要找到一條實際的路徑，
只檢查每個頂點「有一個鄰居讓等號成立」還不夠，下面會給反例。 假設 $d[v]$ 允許是 $\infty$（宣稱走不到）。

\begin{pseudo}{檢查 $d$ 是不是最短距離}
\Function{Verify-Distances}{$G, w, s, d$}
    \If{$d[s] \ne 0$}
        \State \Return \FALSE
    \EndIf
    \For{每條邊 $\{u, v\} \in E$}
        \If{$d[v] > d[u] + w(u, v)$ \OR $d[u] > d[v] + w(u, v)$}
            \State \Return \FALSE \Comment{有更短的走法}
        \EndIf
    \EndFor
    \For{每個頂點 $v$}
        \State $\mathit{seen}[v] \gets \FALSE$
    \EndFor
    \State $\mathit{seen}[s] \gets \TRUE$；$Q \gets \langle s \rangle$
    \While{$Q$ 不是空的}
        \State $u \gets$ \Call{Dequeue}{$Q$}
        \For{每個 $v \in \mathit{Adj}[u]$}
            \If{$\mathit{seen}[v] = \FALSE$ \AND $d[v] = d[u] + w(u, v)$}
                \State $\mathit{seen}[v] \gets \TRUE$；\Call{Enqueue}{$Q, v$}
            \EndIf
        \EndFor
    \EndWhile
    \For{每個頂點 $v$}
        \If{$d[v] < \infty$ \AND $\mathit{seen}[v] = \FALSE$}
            \State \Return \FALSE \Comment{沒有路徑能達到 $d[v]$}
        \EndIf
    \EndFor
    \State \Return \TRUE
\EndFunction
\end{pseudo}

第 4--6 行檢查每條邊的兩個方向（無向邊可以兩個方向走）。 約定 $\infty \le \infty$ 成立、
$\infty + w = \infty$，所以兩端都是 $\infty$ 的邊不會被判錯，一端有限一端 $\infty$ 的邊會被第 5 行抓到。
第 9--14 行是資結書第~21~章的 BFS，只走「等號成立」的邊（叫它\textbf{緊的邊}），$d[u] = \infty$ 的頂點
不會被走到。

\textbf{正確性}。 證明回傳 \TRUE 若且唯若每個 $v$ 都有 $d[v] = \delta(s, v)$。

$(\Leftarrow)$ 設 $d = \delta$。 $d[s] = \delta(s, s) = 0$，第 2 行通過。 由三角不等式（資結書引理~25.3，
無向邊兩個方向各是一條有向邊），第 5 行都不成立。 設 $v$ 有 $d[v] < \infty$，取一條從 $s$ 到 $v$ 的最短
路徑 $\langle s = u_0, u_1, \ldots, u_k = v \rangle$。 由資結書引理~25.2，每一段前綴都是最短路徑，所以
$\delta(s, u_t) = \delta(s, u_{t-1}) + w(u_{t-1}, u_t)$，路徑上每條邊都是緊的。 BFS 沿著緊的邊從 $s$ 出發，
一定會走到 $v$：對 $t$ 歸納，$u_{t-1}$ 被標記後會被取出，檢查它的鄰居 $u_t$ 時，$u_t$ 若還沒被標記就會
被標記。 所以第 16 行的條件都不成立，回傳 \TRUE。

$(\Rightarrow)$ 設回傳 \TRUE。
\begin{itemize}
\item \textbf{$d[v] \le \delta(s, v)$}。 $\delta(s, v) = \infty$ 時顯然。 否則取一條最短路徑
$\langle u_0 = s, \ldots, u_k = v \rangle$，第 5 行不成立表示 $d[u_t] \le d[u_{t-1}] + w(u_{t-1}, u_t)$，
從 $d[u_0] = 0$ 一路加上去，$d[v] \le w(\text{路徑}) = \delta(s, v)$。
\item \textbf{$d[v] \ge \delta(s, v)$}。 $d[v] = \infty$ 時顯然。 否則第 15--17 行沒有回傳 \FALSE，$v$ 被
BFS 標記過，沿著 BFS 的前驅從 $v$ 走回 $s$，是一條每條邊都緊的路徑 $\langle s = u_0, \ldots, u_k = v \rangle$。
把等式 $d[u_t] = d[u_{t-1}] + w(u_{t-1}, u_t)$ 加起來，$d[v] = d[s] + w(\text{路徑}) = w(\text{路徑})
\ge \delta(s, v)$。
\end{itemize}
兩邊合起來 $d[v] = \delta(s, v)$。

\textbf{為什麼不能只檢查「每個頂點有一條緊的邊」}。 權重可以是 $0$。 反例：邊 $\{s, u\}$ 權重 $10$、
$\{u, v\}$ 權重 $0$，宣稱 $d = (d[s], d[u], d[v]) = (0, 5, 5)$。 每條邊的三角不等式都成立
（$5 \le 0 + 10$、$0 \le 5 + 10$、$5 \le 5 + 0$），而且 $u$、$v$ 各有一條緊的邊（彼此之間那條）。 但
真正的 $\delta(s, u) = \delta(s, v) = 10$。 上面的演算法從 $s$ 出發，$s$--$u$ 不緊（$0 + 10 \ne 5$），
走不到 $u$、$v$，第 17 行回傳 \FALSE。 權重都是正數時，緊的邊一定從 $d$ 小的指向 $d$ 大的，沿著它往回
走 $d$ 嚴格變小，最後只能停在 $s$，「每個頂點有一條緊的邊」就夠了；題目允許 $0$，所以要做 BFS。

\textbf{時間}。 第 2 行 $O(1)$；第 4--6 行每條邊一次，$O(m)$；BFS 用鄰接串列，每個頂點進出佇列最多一次、
每條邊在兩端各看一次，$O(n + m)$；最後掃一遍頂點 $O(n)$。 合計 $O(n + m) = O(|V| + |E|)$，是線性時間。
沒有用到任何優先佇列，比重新跑一次 Dijkstra（$O(m \log n)$）再逐一比對快。
\end{pastsoln}
